\chapter{Sistema fisico e domanda del progetto}
\section{Manipolatore e attuatori}
La traccia assegna un manipolatore planare 2R con attuazione elastica in
serie. Le coordinate dei link sono $q=[q_1,q_2]^T$, quelle dei rotori
$\theta=[\theta_1,\theta_2]^T$. La deformazione delle molle è $\theta-q$.
Le equazioni, con disturbo fisico di coppia $d$ sui link, sono
\begin{align}
M(q)\ddot q+C(q,\dot q)\dot q+G(q)&=K_s(\theta-q)+d,\\
J_m\ddot\theta+B_m\dot\theta&=\tau-K_s(\theta-q).
\end{align}
Con masse concentrate alle estremità,
\begin{align*}
M_{11}&=(m_1+m_2)l_1^2+m_2l_2^2+2m_2l_1l_2\cos q_2,\\
M_{12}=M_{21}&=m_2l_2^2+m_2l_1l_2\cos q_2,\qquad M_{22}=m_2l_2^2,\\
C&=m_2l_1l_2\sin q_2\begin{bmatrix}-\dot q_2&-\dot q_1-\dot q_2\\\dot q_1&0\end{bmatrix},\\
G&=\begin{bmatrix}(m_1+m_2)gl_1\cos q_1+m_2gl_2\cos(q_1+q_2)\\m_2gl_2\cos(q_1+q_2)\end{bmatrix}.
\end{align*}
Il comando $u$ passa attraverso la dinamica degli azionamenti prima di
diventare coppia motore $\tau$. Questa dinamica è distinta da inerzia e
smorzamento dei rotori, già presenti nelle equazioni meccaniche.
\scheme{sea_physical}{Sistema SEA. La coppia elastica agisce con segni opposti su link e rotori; il disturbo entra sui link.}
\begin{table}[htbp]\centering
\begin{tabular}{llll}\toprule
Parametro&Valore&Parametro&Valore\\\midrule
$m_1,m_2$&$1,\ 0.8$ kg&$l_1,l_2$&$0.4,\ 0.3$ m\\
$J_{m1},J_{m2}$&$0.01$ kg m$^2$&$B_{m1},B_{m2}$&$0.02$ N m s/rad\\
$k_1,k_2$&$50,\ 40$ N m/rad&$g$&$9.81$ m/s$^2$\\
$\omega_{n0},\zeta_0$&$40$ rad/s, $0.8$&$\tau_0$&$0.015$ s\\\bottomrule
\end{tabular}
\caption{Meccanica della traccia e scelte nominali degli attuatori già adottate.}
\end{table}
\section{Equilibrio e coordinate}
Manteniamo $\bar q=[\pi/2,0]^T$, con velocità nulle,
$\bar\tau=G(\bar q)$ e $\bar\theta=\bar q+K_s^{-1}G(\bar q)$.
In questa configurazione $G(\bar q)=0$ analiticamente, quindi
$\bar\theta=\bar q$: piccoli residui numerici del coseno non costituiscono
una deformazione fisica. Lo stato meccanico è
$x_m=[q;\dot q;\theta;\dot\theta]$. La linearizzazione usa deviazioni
rispetto a $\bar x_m$ e ha otto stati; gli attuatori nominali ne aggiungono sei.
L'uscita di controllo robusto è $q$, mentre LQG/LQGI usano quattro encoder
$y=[q;\theta]$ con velocità stimate. Non modifichiamo queste architetture.
\section{Che cosa vogliamo confrontare}
La domanda centrale è se conservare la struttura fisica dell'incertezza
permetta di ottenere prestazioni robuste diverse da quelle ottenibili con
una descrizione concentrata. Il confronto richiede lo stesso nominale, gli
stessi intervalli fisici, gli stessi disturbi e pesi. Una differenza dei
margini può dipendere sia dall'inviluppo usato sia dagli accoppiamenti
aggiunti: nel seguito separiamo questi effetti.

\chapter{Incertezze meccaniche e costruzione manuale della LFT}
\section{Scelte e normalizzazione}
Manteniamo incerti soltanto $k_1,k_2,J_{m1}$, con variazione indipendente
del $\pm10\%$. Le rigidezze spostano i modi elastici; l'inerzia riflessa
modifica la dinamica del primo rotore. Gli intervalli sono ipotesi
ingegneristiche, non tolleranze identificate. Masse, lunghezze, $J_{m2}$,
smorzamenti e gravità restano noti nel modello. Aggiungere altri parametri
senza dati renderebbe meno interpretabile il confronto.
Per un parametro simmetrico scriviamo
\[
p=p_0(1+r_p\delta_p),\qquad \delta_p\in[-1,1].
\]
MATLAB usa \texttt{ureal} per il parametro fisico. L'oggetto normalizzato
$\delta_p$ è invece adimensionale. Le occorrenze dello stesso parametro
non sono perturbazioni indipendenti: devono utilizzare lo stesso scalare.
\section{Convenzione LFT e segnali}
Adottiamo l'interconnessione superiore
\[
\begin{bmatrix}z\\y\end{bmatrix}=
\begin{bmatrix}M_{11}&M_{12}\\M_{21}&M_{22}\end{bmatrix}
\begin{bmatrix}w\\u\end{bmatrix},\quad w=\Delta z,
\]
da cui
\[
F_u(M,\Delta)=M_{22}+M_{21}\Delta(I-M_{11}\Delta)^{-1}M_{12}.
\]
Con questo ordinamento MATLAB ricompone il modello con
\texttt{lft(Delta,M)}. Il comando \texttt{lftdata} restituisce una coppia
ricomponibile in questo modo. Per la sintesi, invece, $F_l(P,K)$ chiude
il controllore sulle ultime porte della pianta generalizzata: i due usi
della LFT hanno segnali e significati differenti.
\scheme{manual_mechanics}{LFT superiore della meccanica: $z$ è prodotto dalla parte nominale e $w$ è la forza restituita dall'incertezza.}
\section{Derivazione dalle forze elastiche}
Poniamo $e_s=\Delta\theta-\Delta q$ e introduciamo
\[
z_{k_i}=r_k k_{i0}e_{s,i},\quad w_{k_i}=\delta_{k_i}z_{k_i}.
\]
La coppia elastica è $K_{s0}e_s+w_k$. Lo stesso $w_k$ entra con segno
positivo nell'accelerazione dei link e negativo in quella dei rotori.
Non occorrono due incertezze indipendenti per descrivere i due effetti.
Per l'inerzia del primo motore spostiamo la parte incerta a destra:
\[
J_{10}\Delta\ddot\theta_1=\tau_1-\tau_{s,1}-B_{10}\Delta\dot\theta_1+w_J,
\quad z_J=-r_JJ_{10}\Delta\ddot\theta_1,\quad w_J=\delta_Jz_J.
\]
Il segnale $z_J$ contiene anche un termine diretto in $w_J$. Questo
feedback algebrico realizza la divisione per $1+r_J\delta_J$;
è ben posto perché tale quantità è almeno $0.9$.
Se $A_0,B_0$ sono le matrici nominali, la costruzione è
\[
\dot x=A_0x+B_0\tau+B_ww,\qquad
B_w=\begin{bmatrix}0&0\\M_0^{-1}&0\\0&0\\-J_0^{-1}&J_0^{-1}e_1\end{bmatrix}.
\]
La prima colonna a blocchi comprende i due canali elastici e l'ultima
il canale inerziale. Le righe di $z$ seguono direttamente le equazioni sopra.
\texttt{build\_manual\_mechanical\_lft.m} costruisce $M_m$ da queste matrici;
nessuna decomposizione automatica è usata per costruirlo.
\section{Verifica riproducibile}
I test impongono gli stessi valori fisici a modello manuale e modello
\texttt{ureal}, inclusi nominale, estremi e un punto interno. Confrontano
le risposte frequenziali, con tolleranza assoluta $10^{-7}$, e verificano
la ricomposizione automatica. Non richiediamo matrici $M$ uguali:
realizzazioni e numero di ripetizioni possono essere diversi pur descrivendo
la stessa mappa ingresso--uscita. Gli esiti sono salvati nella cartella
della revisione.
\IfFileExists{generated/salaris_blocks.tex}{\begin{table}[htbp]\centering\small
\begin{tabular}{lllll}\toprule
Costruzione & Parametro & Tipo & Dim. elem. & Occorrenze\\\midrule
Meccanica automatica & Jm1 & ureal & 1 & 1\\
Meccanica automatica & k1 & ureal & 1 & 1\\
Meccanica automatica & k2 & ureal & 1 & 1\\
Meccanica manuale & dJm1 & ureal & 1 & 1\\
Meccanica manuale & dk1 & ureal & 1 & 1\\
Meccanica manuale & dk2 & ureal & 1 & 1\\
Attuatore automatico & Dres1 & ultidyn & 1 & 1\\
Attuatore automatico & tau1 & ureal & 1 & 1\\
Attuatore automatico & wn1 & ureal & 1 & 5\\
Attuatore automatico & zeta1 & ureal & 1 & 1\\
Attuatore manuale & Dres1 & ultidyn & 1 & 1\\
Attuatore manuale & dtau1 & ureal & 1 & 1\\
Attuatore manuale & dwn1 & ureal & 1 & 2\\
Attuatore manuale & dzeta1 & ureal & 1 & 1\\
\bottomrule\end{tabular}
\caption{Strutture estratte con lftdata: stesso ingresso--uscita, realizzazioni diverse. Le occorrenze di un parametro condividono lo stesso scalare.}
\end{table}
}{}

\chapter{Due rappresentazioni degli attuatori}
\section{Nominale e parametri fisici}
Per ciascun canale definiamo
\[
H(s;\omega_n,\zeta)=\frac{\omega_n^2}{s^2+2\zeta\omega_ns+\omega_n^2},\quad
D_p(s;\tau)=\frac{1-\tau s/2}{1+\tau s/2}.
\]
Le ipotesi precedenti sono conservate: $\omega_n\in[36,44]$ rad/s,
$\zeta\in[0.72,0.88]$, $\tau\in[0.012,0.018]$ s. I due attuatori
hanno lo stesso nominale e parametri indipendenti; non introduciamo una
correlazione comune non documentata.
\section{Il residuo di Padé non è il parametro del ritardo}
Padé1 ha uno zero in $+2/\tau$. Di conseguenza
$e^{-\tau s}/D_p-1$ non è in generale una perturbazione stabile.
Un grafico del solo modulo non risolve questo problema. La precedente
famiglia moltiplicativa M4 resta una famiglia razionale legittima;
l'inclusione automatica di tutti i ritardi esatti non segue da quel grafico.
Usiamo invece il residuo additivo al fattore di ritardo:
\[
D=e^{-\tau s}=D_p(\tau)+W_{res}\Delta_{res},\quad
W_{res}=2f\left(\frac{\tau_{max}s}{0.5+\tau_{max}s}\right)^3,
\qquad f=1.08.
\]
$\delta_\tau$ è reale e costante; $\Delta_{res}$ è un sistema dinamico
stabile con norma al più uno. Per verificare il peso poniamo $x=\omega\tau$.
Dalla fase di Padé si ricava
\[
|e^{-jx}-D_p(j\omega;\tau)|
=2\left|\sin\frac{x-2\arctan(x/2)}2\right|
\le\min(2,x^3/12).
\]
Il peso domina questa maggiorazione per $\tau\le\tau_{max}$: il rapporto
massimo dei due rami si trova in $x=\sqrt[3]{24}$ ed è minore di uno.
I tre zeri del peso nell'origine sono cancellati dai tre zeri del residuo
esatto; non compaiono poli instabili nella perturbazione normalizzata.
La famiglia dinamica ammette comunque residui che non corrispondono ad alcun
ritardo fisico: è un ampliamento da dichiarare anche per B.
\section{Modello B: parametri separati e LFT manuale}
Il modello parametrico è
\[
G_{Bi}=H(\omega_i,\zeta_i)\,[D_p(\tau_i)+W_{res}\Delta_{res,i}].
\]
La costruzione manuale usa tre stati ausiliari $x,v,h$:
\begin{align*}
e&=u-x,& t&=\omega_0e+r_\omega\omega_0w_{\omega1},\\
z_{\omega1}&=e,& z_{\omega2}&=t-2(\zeta_0v+r_\zeta\zeta_0w_\zeta),\\
\dot x&=v,&\dot v&=\omega_0z_{\omega2}+r_\omega\omega_0w_{\omega2},\\
z_\zeta&=v,&\dot h&=(x-h)/(\tau_0/2)-r_\tau w_\tau,\\
z_\tau&=\dot h,&z_{res}&=x,\\
\tau_{out}&=2h-x+W_{res}w_{res}.&&
\end{align*}
Chiudiamo $w_{\omega1}=\delta_\omega z_{\omega1}$ e
$w_{\omega2}=\delta_\omega z_{\omega2}$ con lo \emph{stesso} parametro.
Inoltre $w_\zeta=\delta_\zeta z_\zeta$, $w_\tau=\delta_\tau z_\tau$ e
$w_{res}=\Delta_{res}z_{res}$. Pertanto
\[
\Delta_i=\operatorname{diag}(\delta_{\omega i}I_2,\delta_{\zeta i},\delta_{\tau i},\Delta_{res,i}).
\]
La ripetizione di $\delta_\omega$ conserva la relazione fra numeratore,
denominatore e smorzamento. Rendere indipendenti le due occorrenze
cambierebbe l'insieme fisico. Il peso residuo aggiunge tre stati nominali
per canale; non cambia la risposta nominale esterna per $\Delta_{res}=0$.
La costruzione è ripetuta per entrambi i canali SISO senza \texttt{lftdata}
nel percorso manuale. Una seconda implementazione automatica serve soltanto
al confronto, verificato anche con un residuo dinamico non nullo.
\scheme{actuator_parametric}{LFT manuale di ciascun attuatore B. Cinque segnali di interconnessione, tre parametri reali indipendenti e un residuo dinamico.}
\section{Modello A: un unico blocco pieno}
Scegliamo
\[
G_A=H_0\,[D_{p0}I_2+W_A\Delta_A],\qquad
\Delta_A\in\mathcal{H}_\infty^{2\times2},\quad\|\Delta_A\|_\infty\le1.
\]
``Un unico blocco'' significa un blocco complesso pieno $2\times2$,
non uno scalare comune e non una diagonale. Può descrivere accoppiamenti
incrociati assenti nei due attuatori SISO. La perturbazione è additiva
rispetto al fattore Padé, con $H_0$ comune: non si inverte lo zero instabile
di Padé. $H_0$ mantiene il decadimento ad alta frequenza.
\scheme{actuator_concentrated}{Attuatori concentrati. La variante diagonale cambia soltanto la struttura di $\Delta_A$, mantenendo il peso identico.}
\section{Inclusione degli insiemi e costo dell'inviluppo}
La copertura non è affidata soltanto a un fit di campioni. Con
$a=1/\tau_{max}$ scegliamo $W_A=C_A s/(s+a)$. Ponendo
$R=H/H_0$, $\alpha=\omega_n^2/\omega_0^2-1$ e
$\beta=2\zeta_0\omega_n^2/\omega_0-2\zeta\omega_n$, otteniamo
\[
R-1=\frac{\alpha s^2+\beta s}{s^2+2\zeta\omega_ns+\omega_n^2}.
\]
La disuguaglianza
$|s^2+2\zeta\omega_ns+\omega_n^2|_{s=j\omega}
\ge\zeta_{min}(\omega^2+\omega_n^2)$ permette le maggiorazioni
\begin{align*}
B_R&=1+\frac{\alpha_{max}+\beta_{max}/(2\omega_{min})}{\zeta_{min}},\\
B_{QR}&=\frac{\sqrt{\alpha_{max}^2+(\beta_{max}/\omega_{min})^2}}{\zeta_{min}}
\max(1,a/\omega_{min}),\\
B_{QD}&=\max|\tau-\tau_0|\max(a,2/\tau_{min}),\\
C_A&=B_{QR}+B_{QD}+2fB_R\sqrt{12/(2.25)^3}.
\end{align*}
Infatti $Q_R=(R-1)(s+a)/s$ è stabile e proprio; il termine
$(D_p-D_{p0})(s+a)/s$ ha norma al più $B_{QD}$.
Il rapporto $W_{res}(s+a)/s$ ha norma $2f\sqrt{12/(2.25)^3}$.
La disuguaglianza triangolare dà così $\|\Delta_{Ai}\|_\infty\le1$
per ogni membro di B. Segue
\[
\mathcal B\subseteq\mathcal A_{diag}\subseteq\mathcal A_{full}.
\]
La costruzione è verificabile, ma la disuguaglianza triangolare può essere
larga. I campioni mostrano quanto il peso sovrastimi l'inviluppo; non sono
usati per proclamare una prova globale. B--A diagonale misura insieme
perdita di struttura fisica e maggiorazione del peso. A diagonale--A pieno
isola gli accoppiamenti, perché il peso è esattamente lo stesso.
\resultfigure{coverage}{Peso analitico A, inviluppo campionato di B e peso del residuo.}
\section{Interconnessione con il manipolatore}
Le incertezze meccaniche sono identiche nelle tre formulazioni. Gli ingressi
fisici sono $[u;d]$; il disturbo non attraversa l'attuatore. La composizione
mantiene separati $\Delta_m$ e $\Delta_a$. La verifica del nominale confronta
direttamente le risposte di tutti e tre i modelli.
\scheme{combined_uncertainty}{Modello complessivo: cambia la rappresentazione degli attuatori, mentre meccanica e porte fisiche restano comuni.}

\chapter{Controllori e confronto robusto}
\section{Specifiche comuni e controllori conservati}
I target temporali M5 restano: assestamento entro 0.8 s, overshoot entro
10\%, errore statico relativo entro 5\% e cross-coupling entro 0.25
(quest'ultimo è un obiettivo preliminare). La scala di coppia 12 N m
non è una saturazione hardware. Non cambiamo questi target per far passare
un candidato e distinguiamo il loro esito da quello della norma pesata.
Manteniamo la banca sviluppata: LQG, LQGI, $H_\infty$ a uno e due gradi
di libertà, PID ottenuto con \texttt{hinfstruct} e $K_\mu$. La formulazione
a un grado di libertà deriva da sensibilità mista; il progetto esplicito
a due gradi di libertà è ottenuto con \texttt{hinfsyn}. Gli script di
sintesi e le matrici dei controllori sono conservati.
Per un confronto omogeneo valutiamo per tutti
\[
w=[r;d],\quad z_p=[W_e(q-M_or);W_uu],\quad u=K_rr+K_yy.
\]
Usiamo i pesi del controllore nominale finale. Questa NP comune non coincide
necessariamente con il costo proprio di LQG o la precedente sensibilità
mista a un grado di libertà. Un fallimento di questa NP non invalida
le verifiche originarie su un obiettivo diverso.
\IfFileExists{generated/salaris_nominal.tex}{\begin{table}[htbp]\centering\small
\begin{tabular}{llllll}\toprule
K & Ordine & $T_s$ [s] & OS [\%] & Errore DC & Picco $u$ [N m]\\\midrule
LQG & 14 & $0.675$ & $0.411$ & $1.39\cdot10^{-14}$ & $4.75$\\
LQGI & 16 & $5.05$ & $0$ & $1.11\cdot10^{-16}$ & $0.822$\\
Hinf1 & 18 & $1.56$ & $125$ & $0.142$ & $13.9$\\
Hinf2 & 22 & $1.46$ & $0$ & $0.00517$ & $0.868$\\
PID & 4 & $2.98$ & $11$ & $0$ & $3.88$\\
Kmu & 78 & $1.52$ & $0$ & $0.06$ & $0.798$\\
\bottomrule\end{tabular}
\caption{Risultati nominali deterministici: massimo sui due gradini di riferimento di 5 gradi. Assestamento al 2 per cento attorno al valore DC effettivo.}
\end{table}
}{}
\scheme{synthesis_analysis}{Pianta generalizzata e controllore. Il circuito superiore descrive l'incertezza; quello inferiore chiude il controllo. $z_p$ è distinto da $z_\Delta$.}
\section{NS, NP, RS e RP}
NS richiede stabilità interna del nominale; NP richiede inoltre
$\|T_{wz_p}(0)\|_\infty<1$, dove $(0)$ indica incertezze nulle e non
frequenza nulla. RS riguarda tutte le perturbazioni ammesse; RP richiede
anche la prestazione pesata per tutte le perturbazioni.
Un limite inferiore del margine maggiore di uno significa PASS; un limite
superiore minore di uno significa FAIL; un intervallo che contiene uno
resta inconclusivo. RP implica RS. Se RS fallisce prima del raggio uno,
usiamo $[0,\mathrm{RS}_{upper}]$ come sola maggiorazione del margine RP.
Non accreditiamo intervalli numericamente invertiti di \texttt{robgain}.
Se NP fallisce al nominale, il margine RP è zero.
\IfFileExists{generated/salaris_comparison.tex}{\begin{table}[htbp]\centering\small
\begin{tabular}{lllllll}\toprule
K & Modello & NS & NP & RS & RP $\le$ & Banda [rad/s]\\\midrule
LQG & A pieno & 1 & 0 & $[0.148,0.148]$ & $0$ & $5.4$\\
LQG & A diag. & 1 & 0 & $[0.149,0.15]$ & $0$ & $5.4$\\
LQG & B & 1 & 0 & $[0.261,0.262]$ & $0$ & $5.4$\\
LQGI & A pieno & 1 & 0 & $[0.143,0.143]$ & $0$ & $0.812$\\
LQGI & A diag. & 1 & 0 & $[0.145,0.145]$ & $0$ & $0.812$\\
LQGI & B & 1 & 0 & $[0.253,0.254]$ & $0$ & $0.812$\\
Hinf1 & A pieno & 1 & 0 & $[0.0731,0.0733]$ & $0$ & $21.4$\\
Hinf1 & A diag. & 1 & 0 & $[0.0758,0.0763]$ & $0$ & $21.4$\\
Hinf1 & B & 1 & 0 & $[0.0973,0.0999]$ & $0$ & $21.4$\\
Hinf2 & A pieno & 1 & 1 & n.i. & n.i. & $2.64$\\
Hinf2 & A diag. & 1 & 1 & n.i. & n.i. & $2.64$\\
Hinf2 & B & 1 & 1 & $[0.0116,0.0116]$ & $0.0116$ & $2.64$\\
PID & A pieno & 1 & 1 & $[0.0135,0.0135]$ & $0.0135$ & $1.12$\\
PID & A diag. & 1 & 1 & $[0.0229,0.023]$ & $0.023$ & $1.12$\\
PID & B & 1 & 1 & $[0.07,0.0706]$ & $0.0706$ & $1.12$\\
Kmu & A pieno & 1 & 1 & $[0.35,0.351]$ & $0.351$ & $2.39$\\
Kmu & A diag. & 1 & 1 & $[0.354,0.354]$ & $0.354$ & $2.39$\\
Kmu & B & 1 & 1 & $[0.88,0.883]$ & $0.883$ & $2.39$\\
\bottomrule\end{tabular}
\caption{Stessi controllori su tre famiglie. RP è maggiorato da RS quando RS fallisce. NP usa la stessa mappa pesata per tutti.}
\end{table}
}{Analisi comparativa NON VALIDATA.}
\resultfigure{rs_comparison}{Margini inferiori RS per gli stessi controllori. La soglia è uno.}
\IfFileExists{generated/salaris_findings.tex}{Per il controllore consolidato $K_\mu$, i margini RS inferiori sono A pieno: $0.35$, A diag.: $0.354$, B: $0.88$. Questi valori riguardano le nuove famiglie, non sostituiscono i margini storici di M4. Le frequenze critiche RS corrispondenti sono A pieno: $13.3$ rad/s, A diag.: $13.1$ rad/s, B: $63.6$ rad/s.
 I margini indicati come n.i. violano numericamente la monotonia imposta dalle inclusioni: restano quantitativamente inconclusivi. I valori grezzi sono conservati nel CSV; non si usano per sostenere un vantaggio fra modelli.
}{}
\section{Sintesi sulle due rappresentazioni}
Una sintesi $H_\infty$ ordinaria usa il nominale comune: non avrebbe senso
presentare due controllori diversi soltanto cambiando l'etichetta
dell'incertezza. Per ogni livello eseguiamo una sintesi nominale comune
e due D--K, una su A pieno e una su B. Ogni candidato viene analizzato
su entrambi i modelli.
La campagna ha un massimo di quattro iterazioni, ordine di fit $D/G$
pari a $[2,1]$, scalature diagonali e seed 20261024. I parametri reali
rimangono reali nell'analisi mista. Sono limiti computazionali, non prove di
ottimalità. Se il problema originale non raggiunge RS/RP, proviamo fattori
$b=0.3$ e poi $b=0.1$, sostituendo $W_e(s)$ e $M_o(s)$ con
$W_e(s/b)$ e $M_o(s/b)$. La richiesta DC resta uguale, mentre la dinamica
richiesta è più lenta. $W_u$ non cambia. La campagna termina quando
entrambe le sintesi superano RS/RP; i candidati non sostituiscono la baseline.
\IfFileExists{generated/salaris_synthesis.tex}{\begin{table}[htbp]\centering\small
\begin{tabular}{lllllll}\toprule
Sintesi & Modello & $b$ & Ordine & RS inf. & RP inf. & Banda [rad/s]\\\midrule
$H_\infty$ & A pieno & $1$ & 22 & $0.00461$ & $0$ & $2.64$\\
$H_\infty$ & B & $1$ & 22 & $0.00889$ & $0$ & $2.64$\\
$\mu$ & A pieno & $1$ & 50 & $0.375$ & $0$ & $0.972$\\
$\mu$ & B & $1$ & 134 & $1.31$ & $0$ & $1.57$\\
$H_\infty$ & A pieno & $0.3$ & 22 & $0.00456$ & $0$ & $0.773$\\
$H_\infty$ & B & $0.3$ & 22 & $0.00886$ & $0$ & $0.773$\\
$\mu$ & A pieno & $0.3$ & 50 & $0.319$ & $0$ & $0.228$\\
$\mu$ & B & $0.3$ & 134 & $0.834$ & $0$ & $0.676$\\
$H_\infty$ & A pieno & $0.1$ & 22 & $0.00473$ & $0$ & $0.257$\\
$H_\infty$ & B & $0.1$ & 22 & $0.00918$ & $0$ & $0.257$\\
$\mu$ & A pieno & $0.1$ & 52 & $0.475$ & $0$ & $0.132$\\
$\mu$ & B & $0.1$ & 132 & $1.54$ & $0$ & $0.155$\\
\bottomrule\end{tabular}
\caption{Sintesi nominale comune e sintesi D--K distinte. Le analisi incrociate complete sono nel CSV synthesis.}
\end{table}
}{Sintesi A/B NON VALIDATA.}
La banda è il minimo dei due canali diagonali di tracking, misurata a
$-3$ dB rispetto al proprio valore DC. Non coincide con il fattore $b$,
né certifica il tempo di assestamento. Una decade di perdita può essere
dichiarata soltanto se emerge dalle bande misurate.
Per i candidati D--K verifichiamo inoltre i livelli comuni
$\|T_{wz_p}\|_\infty<\gamma_{req}$ con $\gamma_{req}=1,3,10$.
Questo equivale a dividere entrambi i pesi per $\gamma_{req}$: non richiede
una nuova ottimizzazione e non cambia la risposta nominale del candidato.
Si rilassa il limite di amplificazione pesata, non si attribuisce una
banda differente al medesimo controllore. RS resta un requisito separato.
\IfFileExists{generated/salaris_relaxation.tex}{\begin{table}[htbp]\centering\small
\begin{tabular}{llllll}\toprule
Modello & $b$ & $\gamma_{req}$ & RS inf. & RP inf. & Banda [rad/s]\\\midrule
A pieno & $1$ & $10$ & $0.375$ & $0$ & $0.972$\\
A pieno & $0.3$ & $10$ & $0.319$ & $0$ & $0.228$\\
A pieno & $0.1$ & $10$ & $0.475$ & $0$ & $0.132$\\
B & $1$ & $3$ & $1.31$ & $1.3$ & $1.57$\\
B & $0.3$ & $10$ & $0.834$ & $0.836$ & $0.676$\\
B & $0.1$ & $3$ & $1.54$ & $1.53$ & $0.155$\\
\bottomrule\end{tabular}
\caption{Primo livello fra 1,3,10 con garanzia RS/RP; se nessuno riesce è mostrato il livello 10, senza attribuirgli una garanzia.}
\end{table}
}{}
\IfFileExists{generated/salaris_synthesis_temporal.tex}{\begin{table}[htbp]\centering\small
\begin{tabular}{llllll}\toprule
K & $b$ & $T_s$ [s] & Errore DC & Picco $u_r$ & Picco $q_d$ [rad]\\\midrule
$\mu$ A & $0.1$ & $\infty$ & $0.246$ & $0.455$ & $1.36$\\
$\mu$ A & $0.3$ & $\infty$ & $0.488$ & $0.383$ & $0.895$\\
$\mu$ A & $1$ & $3.28$ & $0.318$ & $0.56$ & $0.615$\\
$\mu$ B & $0.1$ & $\infty$ & $0.161$ & $0.572$ & $1.01$\\
$\mu$ B & $0.3$ & $6.02$ & $0.121$ & $0.72$ & $0.738$\\
$\mu$ B & $1$ & $2.4$ & $0.377$ & $0.505$ & $0.569$\\
$H_\infty$ & $0.1$ & $\infty$ & $9.62\cdot10^{-5}$ & $0.783$ & $0.0442$\\
$H_\infty$ & $0.3$ & $4.87$ & $0.000229$ & $0.822$ & $0.0424$\\
$H_\infty$ & $1$ & $1.46$ & $0.00519$ & $0.87$ & $0.0457$\\
\bottomrule\end{tabular}
\caption{Prestazioni nominali effettive dei nuovi candidati. Coppie in N m, disturbo lato link 0.5 N m; valori massimi sui due canali.}
\end{table}
}{}

\chapter{Prove rappresentative e integrazioni mirate}
\section{Dieci impianti fisici incerti}
Con seed 20261024 estraiamo dieci vettori di nove parametri normalizzati:
tre meccanici e tre per ciascun attuatore. Simuliamo B con residui dinamici
nulli: è variabilità parametrica Padé, non un campionamento esaustivo di
A o dei residui ammessi da B. Gli stessi impianti sono usati per tutti
i controllori. Un riferimento di 5 gradi e un disturbo lato link di
0.5 N m servono a misurare stabilità osservata, assestamento, errore DC,
picco articolare da disturbo e sforzo. L'orizzonte è 12 s; un assestamento
non osservato è indicato come infinito, non sostituito con 12 s.
\IfFileExists{generated/salaris_samples.tex}{\begin{table}[htbp]\centering\small
\begin{tabular}{llllll}\toprule
K & Stabili & $T_s$ [s] & Errore DC & Picco $q_d$ [rad] & Picco $u_d$ [N m]\\\midrule
LQG & 10/10 & $0.7$ & $0.0363$ & $0.0767$ & $1.36$\\
LQGI & 10/10 & $5.29$ & $2.22\cdot10^{-16}$ & $0.0599$ & $1.38$\\
Hinf1 & 6/10 & $\infty$ & $0.142$ & $0.0753$ & $1.11$\\
Hinf2 & 0/10 & -- & -- & -- & --\\
PID & 4/10 & $2.56$ & $1.11\cdot10^{-16}$ & $0.0999$ & $2.71$\\
Kmu & 10/10 & $1.56$ & $0.06$ & $0.0933$ & $1.2$\\
\bottomrule\end{tabular}
\caption{Dieci impianti parametrici comuni; massimi sui soli campioni stabili.}
\end{table}
}{Simulazioni campionarie NON VALIDATE.}
\section{Dieci condizioni iniziali e confronto lineare/non lineare}
La RAS è esplorata variando lo stato iniziale, non i parametri.
Usiamo lo stesso controllore nominale su lineare e non lineare,
senza riferimenti, disturbi o rumore. Velocità, stati degli attuatori e
del controllore sono inizialmente nulli. Impostiamo $\Delta\theta=\Delta q$:
nel non lineare aggiungiamo $\bar x_m$; nel lineare usiamo le deviazioni.
Le dieci coppie, in gradi, sono
\[
(0,0),(5,0),(15,0),(30,0),(0,5),(0,15),(0,30),
(15,15),(-15,-15),(30,-30).
\]
Nell'ultimo secondo di una simulazione di 12 s richiediamo posizioni
entro 0.5 gradi e velocità entro 0.02 rad/s. Separatamente misuriamo
\[
J(x_0)=\max_{t\in[0,12]}\|q_{NL}(t;x_0)-q_{LIN}(t;x_0)\|_\infty,
\quad \varepsilon=0.5^\circ.
\]
La soglia è il 10\% del riferimento di prova: indica una discrepanza
transitoria contenuta rispetto al movimento richiesto. Un campione può
convergere senza soddisfare $J\le\varepsilon$. Un errore o timeout del
solver è un esito numerico ignoto, distinto dalla mancata convergenza
in un'integrazione completata. Nessuna funzione di Lyapunov è necessaria.
\IfFileExists{generated/salaris_ras.tex}{\begin{table}[htbp]\centering\small
\begin{tabular}{lllll}\toprule
K & Convergenti & Convergenti, $J\le\varepsilon$ & Ignote & Max $J$ [rad]\\\midrule
LQG & 10/10 & 6/10 & 0 & $0.145$\\
LQGI & 10/10 & 6/10 & 0 & $0.177$\\
Hinf1 & 10/10 & 3/10 & 0 & $0.142$\\
Hinf2 & 9/10 & 2/10 & 1 & $0.809$\\
PID & 3/10 & 2/10 & 2 & $2.23$\\
Kmu & 10/10 & 5/10 & 0 & $0.0738$\\
\bottomrule\end{tabular}
\caption{Esplorazione empirica su dieci punti. Gli arresti numerici rimangono ignoti.}
\end{table}
}{RAS e confronto LIN/NL NON VALIDATI.}
\resultfigure{ras_points}{Verde: convergenza e soglia rispettata; arancio: convergenza con scarto eccessivo; croce rossa: non convergente nell'orizzonte; quadrato grigio: esito numerico ignoto. Il segno nero è l'equilibrio in deviazione.}
\resultfigure{linear_nonlinear_error}{Indice $J$ rispetto alla distanza iniziale. Convergenza e fedeltà al lineare sono criteri distinti.}
\resultfigure{paired_responses}{Esempi LIN/NL con $K_\mu$ e condizioni iniziali fisicamente equivalenti.}
\section{Q/R e Kalman}
Mostriamo soltanto tre moltiplicatori di $Q$ ($0.5,1,2$) con $R$ invariata,
in una prova LQR a stato completo che isola il costo dal filtro.
Il prefiltro garantisce lo stesso tracking DC. I grafici sovrapposti
mostrano risposta e comando; non selezioniamo un nuovo controllore.
\resultfigure{qr}{Compromesso risposta--comando al variare del peso sugli stati.}
Per il filtro confrontiamo grandezze statisticamente omogenee.
$Q_n=0.1^2I_4$ descrive i quattro ingressi di coppia di processo;
$R_n=(0.05\pi/180)^2I_4$ descrive i quattro encoder. Sono scale progettuali,
non dati identificati. LQR minimizza $\int(x^TQx+u^TRu)dt$; la matrice
$Q$ deriva dalle scale di posizione, velocità e deformazione già adottate,
mentre $R$ penalizza direttamente la coppia. Kalman stima gli stati non
misurati; LQGI aggiunge due stati che integrano l'errore di posizione.
$Q_n,R_n$ sono intensità di rumore bianco continuo, non varianze di
campioni ZOH. Costruiamo la covarianza stazionaria del sistema congiunto
pianta--stimatore e simuliamo esattamente il processo campionato usando
$\Phi=e^{A\Delta t}$ e $Q_d=P-\Phi P\Phi^T$.
Il periodo è 0.01 s, la durata 600 s e i primi 60 s sono scartati.
La covarianza di $x-\hat x$ è confrontata con Riccati di \texttt{lqe};
la varianza dell'uscita è confrontata con la propria covarianza di Lyapunov.
\IfFileExists{generated/salaris_kalman.tex}{\begin{table}[htbp]\centering\small
\begin{tabular}{lllll}\toprule
Coordinata & Var $q$ teorica & Var $q$ empirica & Var $e_q$ teorica & Var $e_q$ empirica\\\midrule
q1 & $0.00111$ & $0.0011$ & $1.73\cdot10^{-5}$ & $1.71\cdot10^{-5}$\\
q2 & $0.00741$ & $0.00738$ & $6.52\cdot10^{-5}$ & $6.56\cdot10^{-5}$\\
\bottomrule\end{tabular}
\caption{Varianze in rad quadrati. Uscita ed errore di stima sono grandezze diverse; 540 s dopo il transitorio.}
\end{table}
}{Verifica statistica NON VALIDATA.}
LQG con prefiltro è un servo senza modello interno del disturbo costante.
LQGI integra l'errore misurato; $r=0$ dà la corrispondente regolazione.
I robusti a due gradi di libertà ricevono separatamente $r$ e $q$.
Le medesime interfacce sono usate sul non lineare, senza sostituire
lo stato reale alle misure dei controllori.
\section{Riduzione del controllore a 82 stati}
Il controllore MuN1 del precedente studio sul rumore è conservato.
Confrontiamo soltanto riduzioni bilanciate a 70 e 60 stati, riportando
errore relativo $H_\infty$, stabilità, norma pesata, RS/RP, banda,
reiezione del disturbo e amplificazione del rumore encoder sul comando.
Quest'ultima è una norma indotta, non un picco gaussiano simulato.
Nessun ordine inferiore viene promosso automaticamente.
\IfFileExists{generated/salaris_cancellations.tex}{La verifica con minreal a tolleranza $10^{-8}$ passa da 82 a 81 stati, con errore relativo di norma $7.47\cdot10^{-5}$. È una diagnostica delle cancellazioni; il controllore originale rimane conservato.
}{}
\IfFileExists{generated/salaris_reduction.tex}{\begin{table}[htbp]\centering\small
\begin{tabular}{llllll}\toprule
Ordine & Modello & Errore rel. $H_\infty$ & RS inf. & RP inf. & $\|T_{vu}\|_\infty$\\\midrule
82 & A pieno & $1.16\cdot10^{-13}$ & $0.345$ & $0$ & $887$\\
82 & B & $1.16\cdot10^{-13}$ & $0.94$ & $0$ & $887$\\
70 & A pieno & $1.1\cdot10^{-11}$ & $0.345$ & $0$ & $887$\\
70 & B & $1.1\cdot10^{-11}$ & $0.94$ & $0$ & $887$\\
60 & A pieno & $2.42\cdot10^{-8}$ & $0.345$ & $0$ & $887$\\
60 & B & $2.42\cdot10^{-8}$ & $0.94$ & $0$ & $887$\\
\bottomrule\end{tabular}
\caption{Riduzioni del candidato MuN1 a 82 stati. Nessuna promozione automatica.}
\end{table}
}{Riduzione d'ordine NON VALIDATA.}
\IfFileExists{generated/salaris_legacy_reduction.tex}{\begin{table}[htbp]\centering\small
\begin{tabular}{lllll}\toprule
Ordine & NS & Norma nominale & RS inf. & RP inf.\\\midrule
82 & 1 & $0.673$ & $1.85$ & $1.6$\\
70 & 1 & $0.673$ & $1.85$ & $1.6$\\
60 & 1 & $0.673$ & $1.85$ & $1.6$\\
\bottomrule\end{tabular}
\caption{Controllo della riduzione sulla famiglia storica M4, distinta dai nuovi modelli A/B.}
\end{table}
}{}
Verifichiamo anche la famiglia storica M4: conservare i suoi margini è
necessario per confrontare la riduzione con il precedente risultato a
82 stati. Una conferma su M4 non costituisce una garanzia sui nuovi A/B.

\chapter{Conclusioni, limiti e riproducibilità}
B conserva parametri reali, ripetizioni e indipendenza dei due attuatori.
A rende più semplice la struttura di $\Delta$, ampliando l'insieme.
Le inclusioni sono dimostrate per i pesi adottati, ma non ogni differenza
numerica misura soltanto la struttura: anche la maggiorazione del peso
contribuisce alla conservatività. La variante diagonale chiarisce questa distinzione.
Per descrivere i due attuatori SISO con intervalli fisici assegnati, B è
la rappresentazione più informativa. A è utile per studiare anche dinamiche
incrociate non descritte da quei parametri, purché l'ampiezza del peso sia
giustificata. Nel nostro confronto il peso analitico è largo: i risultati
non dimostrano che ogni possibile modello concentrato sia necessariamente
peggiore di ogni modello parametrico.
\IfFileExists{generated/salaris_conclusions.tex}{Con limiti di amplificazione esplicitamente rilassati risultano garantite RS/RP per: B, $b=1$, $\gamma_{req}=3$, banda $1.57$ rad/s; B, $b=0.1$, $\gamma_{req}=3$, banda $0.155$ rad/s. Il livello unitario e i livelli rilassati sono requisiti diversi. La ricerca limitata non dimostra una banda robusta massima né l'impossibilità delle formulazioni prive di un candidato accreditato. Il candidato garantito con banda più alta ha assestamento nominale $2.4$ s ed errore DC relativo $37.7$\%. Questi valori superano i target M5 di 0.8 s e 5\%: la garanzia pesata rilassata non accredita tali target temporali.
}{}
Le garanzie riguardano il modello lineare incerto. Le simulazioni fisiche
e le condizioni iniziali danno evidenza empirica, non certificazioni
della RAS. Le coppie sono non saturate e non sono disponibili limiti
hardware identificati. Gli otto modelli Simulink e gli script accorpati
sono conservati con le precedenti verifiche: non attribuiamo ai nuovi
modelli incerti un collaudo Simulink non effettuato.
Le nuove sintesi A/B sono candidati verificati sul modello lineare. La
RAS e le prove non lineari di questa revisione riguardano la banca
consolidata; la promozione dei nuovi candidati richiede anche il loro
collaudo non lineare e Simulink.
\section{Esecuzione e provenienza dei risultati}
Servono MATLAB R2026b, Control System Toolbox e Robust Control Toolbox;
la regressione simbolica usa Symbolic Math Toolbox. Dalla radice:
\begin{lstlisting}
addpath(genpath('analysis'));
run_salaris_revision('models');
run_salaris_revision('analysis');
run_salaris_revision('synthesis');
run_salaris_revision('simulations');
run_salaris_revision('secondary');
\end{lstlisting}
Le fasi sono sequenziali e salvano in
\texttt{results/data/salaris\_revision/}. I checkpoint evitano di ripetere
sintesi già eseguite; cambiare modelli o specifiche richiede una nuova
cartella dei risultati. \texttt{python tools/build\_salaris\_documents.py}
genera le tabelle dai CSV; il compilatore del repository genera i PDF.
Il branch è una proposta da revisionare e integrare personalmente.
Le trattazioni precedenti e i dati consolidati restano nel repository;
il percorso principale è stato sostituito per ridurre il volume.
I dettagli storici sono nelle altre sezioni di \texttt{report/sections/}
e nella versione precedente in Git.
\section{Riferimenti essenziali}
\emph{Traccia 2 -- Manipolatore SEA}, corso di Controllo Robusto 2026.
Indicazioni del ricevimento del 22 ottobre 2026 riportate nel mandato.
Documentazione MathWorks:
\href{https://www.mathworks.com/help/robust/ref/inputoutputmodel.lftdata.html}{lftdata},
\href{https://www.mathworks.com/help/control/ref/dynamicsystem.pade.html}{Padé}.
